% ============================== ANEXE ==============================
% Anexele nu se numără în limita de pagini a lucrării. Fiecare anexă este
% referită cel puțin o dată în corpul lucrării (vezi Capitolul 3).

\chapter*{Anexa 1 — Capturi de ecran ale interfeței}
\addcontentsline{toc}{chapter}{Anexa 1 — Capturi de ecran ale interfeței}

Această anexă conține capturi de ecran ale aplicației web. Substituenții de mai jos se vor
înlocui cu imaginile reale (\code{\textbackslash includegraphics}) după capturarea lor în
\code{images/}.

\begin{figure}[!htb]
  \capturaplaceholder{Tabloul de bord (dashboard) după autentificare: rezumatul activității și
  punctele de intrare către funcționalități.}
  \caption{Tabloul de bord al utilizatorului.}
\end{figure}

\begin{figure}[!htb]
  \capturaplaceholder{Constructorul de programe: zile, exerciții ordonate cu serii/repetări/pauză
  planificate și rutine atașate.}
  \caption{Constructorul de programe de antrenament.}
\end{figure}

\begin{figure}[!htb]
  \capturaplaceholder{Sesiunea live: înregistrarea seriilor (greutate, repetări, RPE) și adăugarea
  de exerciții suplimentare.}
  \caption{Execuția unei sesiuni de antrenament.}
\end{figure}

\begin{figure}[!htb]
  \capturaplaceholder{Analize de progres: volum în timp, antrenamente pe săptămână, distribuția pe
  grupe musculare și progresia 1RM estimat (grafice Recharts).}
  \caption{Pagina de analize și progres.}
\end{figure}

\begin{figure}[!htb]
  \capturaplaceholder{Marketplace cu căutare: câmp de interogare, fațete (dificultate, split, grupe
  musculare) și evidențierea potrivirilor (\code{<mark>}).}
  \caption{Marketplace-ul de programe cu căutare full-text și fațete.}
\end{figure}

\begin{figure}[!htb]
  \capturaplaceholder{Modul de coaching: căutarea full-text în istoricul unui client, accesibilă
  doar pe baza unei relații active.}
  \caption{Căutarea în istoricul clientului (mod coaching).}
\end{figure}

\begin{figure}[!htb]
  \capturaplaceholder{Panoul analizorului structural: scorul pe 100 de puncte, sub-scorurile și
  avertismentele explicabile cu sugestii.}
  \caption{Rezultatul analizorului structural de programe.}
\end{figure}

% ----------------------------------------------------------------------------
\chapter*{Anexa 2 — Schema bazei de date}
\addcontentsline{toc}{chapter}{Anexa 2 — Schema bazei de date}

Schema este definită integral de migrarea Flyway \code{V1} (28 de tabele), populată cu date
oficiale de \code{V2} și completată cu un index de concurență de \code{V3}. Tabelele, grupate pe
funcționalități, sunt:

\begin{itemize}
  \item \textbf{Autentificare / coaching:} \code{app\_users}, \code{refresh\_tokens},
        \code{coach\_profiles}, \code{coach\_client\_relationships}.
  \item \textbf{Catalog / planificare:} \code{muscle\_groups}, \code{exercises},
        \code{exercise\_muscle\_groups}, \code{routines}, \code{gyms}, \code{equipment}.
  \item \textbf{Programe:} \code{workout\_templates}, \code{template\_stats}, \code{template\_days},
        \code{template\_day\_exercises}, \code{template\_day\_exercise\_muscle\_groups},
        \code{template\_day\_routines}.
  \item \textbf{Sesiuni (istoric imuabil):} \code{workout\_sessions}, \code{session\_exercises},
        \code{session\_exercise\_muscle\_groups}, \code{session\_routines}, \code{workout\_sets}.
  \item \textbf{Marketplace:} \code{template\_votes}, \code{template\_saves},
        \code{template\_use\_events}.
  \item \textbf{Tabele pregătite pentru extinderi (neutilizate încă de cod):}
        \code{coach\_session\_comments}, \code{coach\_template\_assignments},
        \code{template\_analysis\_results}, \code{outbox\_events}.
\end{itemize}

Constrângerile cu rol esențial în integritate sunt: invariantul de vizibilitate publicată a
programelor (\code{CHECK ((visibility='PUBLIC' AND published\_at IS NOT NULL) OR
visibility='PRIVATE')}); anti-autocoaching (\code{CHECK (coach\_user\_id <> client\_user\_id)});
indexul unic parțial pentru un singur antrenament activ per utilizator; indecșii unici parțiali de
unicitate cu \code{WHERE deleted\_at IS NULL}; și cheile externe anulabile (\code{ON DELETE SET
NULL}) către datele sursă, pereche cu coloanele de instantaneu. Definițiile-cheie sunt în Anexa~3.

% ----------------------------------------------------------------------------
\chapter*{Anexa 3 — Listări de cod reprezentative}
\addcontentsline{toc}{chapter}{Anexa 3 — Listări de cod reprezentative}

\noindent\textbf{Listarea 1.} Stilul determinist al analizorului: fiecare sub-scor pornește de la
maxim și scade o valoare fixă pe constatare (extras simplificat din \code{TemplateAnalyzerService}).
\begin{lstlisting}[language=Java]
private int scoreVolumeAndCoverage() {
    int score = 35;
    // O grupa minora nu poate masca o regiune majora absenta:
    if (!lowerBody) { add("NO_LOWER_BODY", HIGH, ...); score -= 12; }
    if (!pulling)   { add("NO_PULL",       HIGH, ...); score -= 12; }
    if (!pushing)   { add("NO_PUSH",       HIGH, ...); score -= 12; }
    for (String bucket : ALL_BUCKETS) {
        // serii saptamanale ponderate: PRIMARY = 1.0, SECONDARY = 0.5
        double volume = weeklyWeightedSets(bucket);
        if      (volume <= 3)  { add("MUSCLE_VOLUME_VERY_LOW", ...);  score -= 3; }
        else if (volume > 25)  { add("MUSCLE_VOLUME_EXCESSIVE", ...); score -= 5; }
        else if (volume > 20)  { add("MUSCLE_VOLUME_HIGH", ...);      score -= 2; }
        // ... alte benzi de volum ...
    }
    return Math.max(0, score);
}

private static String category(int score) {
    if (score >= 80) return "WELL_STRUCTURED";
    return score >= 55 ? "DECENT_STRUCTURE" : "NEEDS_REVIEW";
}
\end{lstlisting}

\noindent\textbf{Listarea 2.} Ierarhia de autorizare (din \code{SecurityConfig}): RBAC pe rute,
peste care modul de coaching adaugă verificarea ReBAC la nivel de serviciu.
\begin{lstlisting}[language=Java]
.authorizeHttpRequests(auth -> auth
    .requestMatchers(POST, "/api/auth/register", "/api/auth/login",
                           "/api/auth/refresh",  "/api/auth/logout").permitAll()
    .requestMatchers(GET,  "/api/health").permitAll()
    .requestMatchers("/api/admin/**").hasRole("ADMIN")
    .requestMatchers("/api/coach/**").hasRole("COACH")
    .requestMatchers("/api/**").authenticated()
    .anyRequest().permitAll())
\end{lstlisting}

\noindent\textbf{Listarea 3.} Definiții SQL cheie: garda de concurență (V3) și tabela de outbox
pregătită din V1 pentru extinderea cu outbox tranzacțional.
\begin{lstlisting}[language=SQL]
-- Un singur antrenament activ per utilizator (migrarea V3):
CREATE UNIQUE INDEX ux_workout_sessions_one_active_per_user
  ON workout_sessions (user_id) WHERE status = 'IN_PROGRESS';

-- Tabela de outbox, definita in V1, neutilizata inca de cod:
CREATE TABLE outbox_events (
  id             uuid PRIMARY KEY DEFAULT gen_random_uuid(),
  aggregate_type varchar(100) NOT NULL,
  aggregate_id   uuid NOT NULL,
  event_type     varchar(100) NOT NULL,
  payload        jsonb NOT NULL DEFAULT '{}'::jsonb,
  status         varchar(30) NOT NULL DEFAULT 'PENDING',
  attempts       int NOT NULL DEFAULT 0,
  available_at   timestamptz NOT NULL DEFAULT now(),
  processed_at   timestamptz,
  created_at     timestamptz NOT NULL DEFAULT now(),
  last_error     text,
  CONSTRAINT check_outbox_status
    CHECK (status IN ('PENDING','PROCESSING','PROCESSED','FAILED'))
);
CREATE INDEX outbox_events_processing_idx
  ON outbox_events (status, available_at, created_at);
\end{lstlisting}

% ----------------------------------------------------------------------------
\chapter*{Anexa 4 — Rezultatele complete ale evaluării}
\addcontentsline{toc}{chapter}{Anexa 4 — Rezultatele complete ale evaluării}

Măsurătorile au fost realizate la 7 dimensiuni ale corpusului (1k–2M), cu 300 de iterații măsurate
(încălzire 30, buget de timp 15\,s; \emph{n} indică numărul de eșantioane când bugetul a limitat
celula). PostgreSQL: \code{tsvector}+GIN ordonat cu \code{ts\_rank\_cd}, fuzzy cu \code{pg\_trgm}.
OpenSearch: aceleași mapări și interogări ca producția.

\begin{table}[!htb]
  \centering
  \begin{tabular}{@{}rrrrr@{}}
  \toprule
  Corpus & Construcție PG & Construcție OS & Dimensiune PG & Dimensiune OS \\
  \midrule
  1.000     & 216\,ms     & 203\,ms     & 1,6\,MB     & 0,3\,MB \\
  10.000    & 1.926\,ms   & 975\,ms     & 14,6\,MB    & 2,9\,MB \\
  50.000    & 7.966\,ms   & 3.481\,ms   & 68,7\,MB    & 14,3\,MB \\
  100.000   & 14.006\,ms  & 6.349\,ms   & 135,5\,MB   & 28,5\,MB \\
  500.000   & 68.752\,ms  & 37.987\,ms  & 664,7\,MB   & 140,8\,MB \\
  1.000.000 & 125.774\,ms & 57.701\,ms  & 1.313,6\,MB & 349,1\,MB \\
  2.000.000 & 239.192\,ms & 114.908\,ms & 2.632,5\,MB & 738,7\,MB \\
  \bottomrule
  \end{tabular}
  \caption{Timpul de construcție a indexului și dimensiunea pe disc, pe dimensiuni de corpus.}
\end{table}

\begin{table}[!htb]
  \centering
  \begin{tabular}{@{}rrr@{}}
  \toprule
  Corpus & PostgreSQL ($p_{50}/p_{95}/p_{99}$) & OpenSearch ($p_{50}/p_{95}/p_{99}$) \\
  \midrule
  1.000     & 1,46 / 2,68 / 3,12        & 3,25 / 5,70 / 10,56 \\
  10.000    & 10,64 / 14,32 / 17,63     & 2,85 / 5,21 / 7,17 \\
  50.000    & 48,70 / 66,59 / 77,97     & 5,17 / 10,05 / 13,98 \\
  100.000   & 100,21 / 117,87 / 136,46  & 6,84 / 10,59 / 12,83 \\
  500.000   & 779,47 / 1021,25 / 1229,16 & 22,75 / 38,09 / 67,26 \\
  1.000.000 & 872,19 / 1199,43 / 1220,09 & 33,64 / 45,24 / 83,90 \\
  2.000.000 & 1863,97 / 2783,12 / 2783,12 & 63,65 / 73,05 / 111,67 \\
  \bottomrule
  \end{tabular}
  \caption{Latența căutării full-text „bench press” (ms), pe dimensiuni de corpus.}
\end{table}

\begin{table}[!htb]
  \centering
  \begin{tabular}{@{}rrr@{}}
  \toprule
  Corpus & PostgreSQL ($p_{50}/p_{95}/p_{99}$) & OpenSearch ($p_{50}/p_{95}/p_{99}$) \\
  \midrule
  1.000     & 14,45 / 18,91 / 21,69          & 2,37 / 3,74 / 5,63 \\
  10.000    & 127,93 / 134,06 / 136,22       & 3,23 / 7,38 / 10,94 \\
  50.000    & 621,74 / 645,63 / 646,08       & 2,66 / 4,95 / 6,59 \\
  100.000   & 1831,47 / 2169,61 / 2169,61    & 3,15 / 5,74 / 7,38 \\
  500.000   & 2447,22 / 8354,43 / 8354,43    & 4,06 / 7,20 / 8,95 \\
  1.000.000 & 7806,25 / 17790,10 / 17790,10  & 3,94 / 9,18 / 13,35 \\
  2.000.000 & 15825,49 / 16399,71 / 16399,71 & 3,89 / 9,38 / 14,14 \\
  \bottomrule
  \end{tabular}
  \caption{Latența căutării fuzzy „benhc” (greșeală de tastare, ms): degradarea \code{pg\_trgm} vs.\
  latența practic constantă a OpenSearch.}
\end{table}

% ----------------------------------------------------------------------------
\chapter*{Anexa 5 — Contractul API (selecție)}
\addcontentsline{toc}{chapter}{Anexa 5 — Contractul API (selecție)}

Toate endpoint-urile, cu excepția celor de autentificare și a stării de sănătate, cer un token JWT
valid. Răspunsurile de eroare folosesc un plic uniform cu un cod stabil (de exemplu
\code{TEMPLATE\_NOT\_FOUND}, \code{ACTIVE\_WORKOUT\_EXISTS}, \code{SEARCH\_UNAVAILABLE}).

\begin{lstlisting}
Autentificare : POST /api/auth/{register,login,refresh,logout};  GET /api/auth/me
Planificare   : /api/exercises;  /api/routines;  /api/gyms;
                /api/gyms/{id}/equipment;  /api/equipment/{id}
Programe      : /api/templates;  /api/templates/{id};
                POST /api/templates/{id}/{publish,unpublish}
Sesiuni       : POST /api/workouts/start;  GET /api/workouts/active;
                POST /api/workouts/{id}/{finish,cancel};
                POST /api/session-exercises/{id}/sets;  PUT|DELETE /api/sets/{id}
Istoric       : GET /api/history;  GET /api/analytics/overview
Marketplace   : GET /api/marketplace/templates;
                POST /api/marketplace/templates/{id}/{vote,save,use}
Analizor      : GET /api/templates/{id}/analysis
Cautare       : GET /api/search/templates;  GET /api/search/workouts;
                POST /api/admin/search/reindex      (doar ADMIN)
Coaching      : /api/coach/clients/... (rol COACH + relatie activa);
                /api/coaching/{invites,coaches}
\end{lstlisting}

% ----------------------------------------------------------------------------
\chapter*{Anexa 6 — Configurație și rulare}
\addcontentsline{toc}{chapter}{Anexa 6 — Configurație și rulare}

Stack-ul rulează prin Docker Compose (PostgreSQL 16, OpenSearch 2.13 cu un volum persistent, și
aplicația). Indicatorul \code{app.search.enabled} controlează activarea căutării.

\begin{lstlisting}
# application.yml (extras)
spring:
  jpa:
    hibernate:
      ddl-auto: validate        # schema detinuta de Flyway
app:
  security:
    jwt:
      issuer: workout-thesis
      access-ttl: 15m
      refresh-ttl: 14d
  search:
    enabled: ${APP_SEARCH_ENABLED:false}
    uri: ${APP_SEARCH_URI:http://localhost:9200}
    templates-index: templates_v1
    sessions-index:  workout_sessions_v1
\end{lstlisting}

\noindent Rularea cu date demonstrative (istoric 2020--prezent) se face cu profilul \code{demo},
care la pornire reconstruiește indexul de căutare:
\begin{lstlisting}[language=bash]
SPRING_PROFILES_ACTIVE=demo docker compose up --build
\end{lstlisting}
